\chapter{Construction~: un modèle de carnet d'ordres, de bout en bout}\label{ch:ml:build-an-order-book-model-end-to-end}

Le dernier modèle du livre est volontairement petit~: une prévision de la seconde suivante du prix milieu, entraînée sur les sessions du simulateur,
validée sans fuite, compilée en quelques centaines de nanosecondes et exploitée dans le même simulateur, où chaque microseconde
qu'il prend est une microseconde de décisions périmées. Chaque étape réutilise un composant construit plus tôt dans le livre, et l'exécution entière forme un seul
graphe reproductible. Le chapitre se termine sur la question que pose le titre de son problème~: combien vaut une microseconde pour ce modèle~?
La réponse n'est pas celle que suggère le titre. Le modèle gagne 21.73~\$ par session de vingt minutes sur six sessions de test, quatre fois
ce que rapporte la même cotation sans lui~; son edge reste inchangé jusqu'à 50 millisecondes d'instant de décision et tombe à zéro vers
106 millisecondes. Sur un marché dont le meilleur cours change quatre fois par seconde, les microsecondes ne valent rien et un dixième de
seconde vaut tout~; ce que coûte la latence dépend de la vitesse du marché, et non de celle du code.

\section{La spécification}

Le modèle, rédigé comme la \omterm{def:ml:the-machine-learning-team:handoff}{spécification de transmission} (handoff) du chapitre~28 le prévoit~:
\begin{itemize}
\item \emph{Marché.} Le simulateur de bourse du livre~10, une plateforme et un instrument à 100~\$ avec un tick d'un cent, des frais par défaut (une
remise de 0.002~\$ par action pour la fourniture de liquidité, 0.003~\$ pour la prise de liquidité), le flux \texttt{firm\_tape} du livre~7 comme flux de fond,
l'agent derrière le modèle de latence du livre~10 (20~microsecondes dans chaque sens).
\item \emph{Décision.} À chaque changement du meilleur cours, l'agent observe~: il calcule dix variables et prévoit la variation du mid
dans la seconde suivante (en ticks), puis agit.
\item \emph{Action.} Coter un lot (100 actions) au meilleur cours de chaque côté que sa position autorise (au plus un lot à l'achat ou à la vente),
sauf du côté que la prévision annonce sur le point d'être balayé~: le cours acheteur est retiré lorsque la prévision est inférieure à $-\theta$, le cours vendeur
lorsqu'elle est supérieure à $\theta$. Une cotation dont le prix reste au meilleur cours est laissée en place, car sa position dans la file mérite d'être conservée. Une
position détenue pendant 30~secondes est clôturée au marché.
\item \emph{Sessions.} Vingt minutes chacune, sur des graines disjointes~: huit pour l'entraînement, trois pour la validation, six pour le test.
\item \emph{Budget et suivi.} L'instant de décision est la latence mesurée du modèle compilé~; les moniteurs d'entrée du chapitre~27 tournent sur
les sessions de test.
\end{itemize}

La stratégie est celle d'un teneur de marché qui utilise le modèle pour éviter de subir la sélection adverse, usage classique d'une prévision à court horizon par un
fournisseur de liquidité (Cartea, Jaimungal et Penalva en traitent la forme optimale). Une version preneuse de liquidité a été essayée en premier puis abandonnée~: sur
ce marché, le mid bouge d'un tick ou plus en une seconde environ cinq fois sur cent, et traverser une fourchette d'un tick en payant
deux fois les frais du preneur coûte 1.6 tick par aller-retour, davantage que ce que vaut ici la moindre prévision.

L'ensemble de l'exécution est un seul graphe \texttt{firm.workflow} (\cref{lst:ml:lob:graph})~: deux étapes de données, une validation croisée purgée, l'ajustement
et son export, le TCN challenger, le choix du seuil sur les sessions de validation, et les sessions de test à chaque latence.
Chaque étape est indexée par son code, ses paramètres, sa graine et ses entrées, si bien qu'un changement en un point quelconque relance exactement ce qui en dépend.

\section{Données, étiquettes et validation}

Les sessions de recherche utilisent la même classe d'agent qu'en production, avec le trading désactivé~: il enregistre son flux sous forme d'événements
du \omterm{def:ml:data-and-feature-stores:store}{feature store} du chapitre~24 (chaque changement du meilleur cours, avec l'incrément de son déséquilibre du flux d'ordres, et chaque transaction avec le signe de son
preneur) et calcule ses variables en ligne à chaque changement du carnet (\cref{lst:ml:lob:onbook}). Les dix variables sont
le déséquilibre au meilleur cours, la fourchette, le déséquilibre du flux d'ordres de Cont, Kukanov et Stoikov sur une et cinq secondes, le
volume négocié signé sur une et cinq secondes, le nombre de transactions sur cinq secondes, la variation du mid sur une et cinq secondes, et le nombre
de mises à jour du carnet sur une seconde. Les mêmes définitions calculées hors ligne sur les événements enregistrés égalent exactement les valeurs en ligne
à chaque décision (le test d'acceptation de \texttt{firm.lobmodel})~; l'entraînement et le service ne peuvent pas diverger, puisqu'ils utilisent le
même code.

L'étiquette est l'\omterm{def:ml:targets-labels-and-sample-weights:fixed}{étiquette à horizon fixe} du chapitre~2, en temps événementiel (\cref{lst:ml:lob:data})~: la variation du mid, en ticks, entre le dernier
événement connu à la décision et celui d'une seconde plus tard, calculée par \texttt{firm.labeling} sur les variations du mid par événement.
La plupart des étiquettes sont nulles (89.4~\%)~: en une seconde, le mid ne bouge généralement pas. Les huit sessions d'entraînement donnent 51\,924 décisions et les
trois sessions de validation 33\,457~; l'activité varie beaucoup d'une session à l'autre.

La validation respecte le temps. Chaque étiquette couvre une seconde, si bien que des lignes voisines partagent leur futur~; la validation croisée purgée à cinq plis du
chapitre~3 (\texttt{firm.cvsplit}, avec un embargo d'une seconde) retire de chaque pli d'entraînement les lignes dont les étiquettes chevauchent
le pli de test. Trois configurations LightGBM de 150 arbres, à 7, 15 et 31 feuilles, obtiennent des coefficients d'information moyens par pli de
0.558, 0.555 et 0.548~; la plus petite l'emporte. Réajustée sur l'ensemble des sessions d'entraînement, elle a un coefficient d'information de 0.625 dans l'échantillon (in-sample)
et de 0.522 sur les sessions de validation.

Le challenger est le \omterm{def:ml:sequence-models-on-order-books:tcn}{réseau convolutif temporel} (TCN) du chapitre~8 (\texttt{firm.lobseq}, 12\,867 paramètres), appliqué à des fenêtres des variables des 16 dernières
décisions, entraîné sur sept sessions d'entraînement et arrêté tôt sur la huitième. Son coefficient d'information de validation est de
0.452. La forêt l'emporte en précision et, comme la section suivante le mesure, c'est aussi le modèle que l'on peut rendre rapide~; Kolm, Turiel et
Westray ont trouvé de même que des réseaux standard sur des entrées stationnaires de flux d'ordres surpassent des réseaux appliqués aux états bruts du carnet.

\section{Le modèle et sa forme compilée}

La forêt retenue compte 150 arbres de six divisions chacun. Exportée sous forme de tableaux plats par \texttt{firm.mlinfer} du chapitre~26, elle reproduit
bit à bit les prédictions de LightGBM sur les lignes de validation~; il en va de même des branches C++20 générées, de la boucle C++ sur les tableaux et du
noyau Rust, sur 200 vecteurs de test communs (\cref{lst:ml:lob:cpp}, \cref{lst:ml:lob:rust}). L'agent lui-même appelle un parcours NumPy des
mêmes tableaux qui additionne les feuilles dans le même ordre (\cref{lst:ml:lob:predict}).

\begin{table}[htbp]
\centering\small
\begin{tabular}{@{}lrrr@{}}
\toprule
chemin de service & médiane & 99\up{e}~centile & 99.9\up{e}~centile\\
\midrule
appel Python de LightGBM (une ligne) & 263 $\mu$s & 560 $\mu$s & 869 $\mu$s\\
passe avant du TCN dans PyTorch (une fenêtre) & 166 $\mu$s & 460 $\mu$s & 718 $\mu$s\\
forêt plate en Python et NumPy & 37 $\mu$s & 98 $\mu$s & 172 $\mu$s\\
forêt plate, boucle C++20 & 1.1 $\mu$s & 2.5 $\mu$s & 9.1 $\mu$s\\
forêt plate, branches C++20 générées & 153 ns & 383 ns & 1.1 $\mu$s\\
\bottomrule
\end{tabular}
\caption{Latence d'une prédiction, mesurée une fois sur un ordinateur portable (Intel Core Ultra 7 155H, WSL2, une machine partagée sans
cœurs isolés~; le coût propre de l'horloge est d'environ 15 nanosecondes). Données~: \texttt{bench\_lobmodel.py},
\texttt{measured\_latency.csv}.}\label{tab:ml:lob:latency}
\end{table}

Le \Cref{tab:ml:lob:latency} donne l'autre côté du budget. Les branches générées répondent en 153 nanosecondes à la médiane, environ
1\,700 fois plus vite que l'appel Python de LightGBM~; le TCN, moins précis, serait aussi mille fois plus lent que la
forêt compilée. Les exécutions de trading utilisent ces mesures comme instant de décision déclaré de l'agent~: 150 nanosecondes pour le modèle compilé,
263 microsecondes pour l'appel de la bibliothèque.

\section{Le trading dans le simulateur}

\begin{definition}[Péremption de la décision]\label{def:ml:build-an-order-book-model-end-to-end:staleness}
La \emph{péremption de la décision}\index{péremption de la décision} d'un ordre est le temps écoulé entre l'événement de marché sur lequel sa décision repose jusqu'à
l'instant où l'ordre quitte le trader~: la latence des données de marché, toute attente pendant que des décisions antérieures sont encore calculées, et
l'instant de décision lui-même. Si les décisions sont déclenchées à un taux $\lambda$ et que chacune prend $L$, l'attente ne reste bornée que tant que
$\lambda L<1$.
\end{definition}

La décision de l'agent (\cref{lst:ml:lob:act}) s'exécute à chaque changement du carnet qu'il observe. Son seuil est choisi sur les sessions de validation
à la latence du modèle compilé~: P\&L net moyen par session de 43.4~\$, 53.9~\$, 48.0~\$ et 50.1~\$ pour $\theta$ de 0.05, 0.1, 0.15 et
0.2 tick, soit $\theta=0.1$. Sur les six sessions de test, à 150 nanosecondes, le modèle gagne 21.73~\$ par session (erreur type
3.90~\$), positif dans les six, à partir de 34.3 entrées par session~; ses frais d'échange nets sont un crédit de 6.73~\$. Le mark-out de ses entrées une seconde
plus tard vaut 0.405 tick en sa faveur. La même cotation sans modèle (qui ne retire jamais de côté) gagne 5.47~\$ (erreur type
2.22~\$) à partir de 58.7 entrées, avec un mark-out d'une seconde de 0.078 tick~: elle est exécutée plus souvent, et surtout par des ordres qui en savent
davantage. La valeur du modèle réside dans les exécutions qu'il refuse.

Les sessions de test passent aussi par le moniteur d'entrées du chapitre~27, chaque session étant comparée aux sessions d'entraînement. Le plus grand \omterm{def:ml:monitoring-and-retraining:psi}{indice de stabilité de la population}
d'une session se situe entre 0.14 et 1.60, très au-dessus de toute règle empirique, et dans chaque session il provient d'une variable
qui compte des événements (mises à jour de la dernière seconde, transactions des cinq dernières secondes), variable laquelle suit l'activité de la session~: les sessions de test
voient entre 2.8 et 6.4 changements du carnet par seconde. Les quatre variables de prix restent entre 0.02 et 0.13. Le modèle est rentable
dans les six sessions. Un moniteur portant sur des variables de comptage doit avoir une référence couvrant toute la gamme d'activité que le modèle rencontrera, ou
des variables normalisées par l'activité, faute de quoi il alerte chaque après-midi calme.

\section{Ce que coûte la latence}

Les sessions de test sont de nouveau négociées, l'instant de décision étant fixé à chaque point d'une grille allant de zéro à 150 millisecondes
(\cref{tab:ml:lob:results}, \cref{fig:ml:lob:pnl}). Tout le reste est identique~: les sessions, le modèle, le seuil.

\begin{table}[htbp]
\centering\small
\begin{tabular}{@{}>{\centering\arraybackslash}p{2.6cm}>{\centering\arraybackslash}p{2.6cm}>{\centering\arraybackslash}p{2.6cm}>{\centering\arraybackslash}p{2.6cm}>{\centering\arraybackslash}p{2.6cm}@{}}
\toprule
instant de décision & péremption médiane & P\&L net par session (\$) & mark-out à 1~s (ticks) & sessions positives\\
\midrule
0 & 0.02 ms & 21.73 & 0.405 & 6 sur 6\\
150 ns & 0.02 ms & 21.73 & 0.405 & 6 sur 6\\
263 $\mu$s & 0.28 ms & 21.73 & 0.405 & 6 sur 6\\
1 ms & 1.02 ms & 21.23 & 0.399 & 6 sur 6\\
10 ms & 10.0 ms & 21.72 & 0.405 & 6 sur 6\\
30 ms & 30.0 ms & 20.63 & 0.394 & 6 sur 6\\
50 ms & 50.0 ms & 20.52 & 0.417 & 6 sur 6\\
70 ms & 70.0 ms & 12.40 & 0.353 & 6 sur 6\\
100 ms & 154 ms & 1.60 & 0.245 & 3 sur 6\\
150 ms & 802 ms & $-12.27$ & 0.148 & 2 sur 6\\
sans modèle & 0.02 ms & 5.47 & 0.078 & \\
\bottomrule
\end{tabular}
\caption{Le modèle sur six sessions de test à chaque instant de décision~: péremption médiane de la décision, P\&L net moyen par session de vingt minutes,
mark-out moyen à une seconde des entrées, et les sessions au P\&L positif~; la dernière ligne donne la cotation sans modèle. Données~:
\texttt{ml\_lobmodel.latency\_table}.}\label{tab:ml:lob:results}
\end{table}

\begin{omfigure}
\begin{tikzpicture}
\begin{axis}[width=11.5cm,height=6cm,xmode=log,xlabel={temps de décision (ms)},ylabel={P\&L net par session (\$)},xmin=1e-4,xmax=300,
  ymin=-22,ymax=30,tick label style={font=\small},label style={font=\small},grid=major,grid style={gray!25},legend pos=south west,
  legend cell align=left,legend style={font=\small},table/col sep=comma]
\addplot[name path=hi,draw=none,forget plot] table[x=ms,y=hi]{figdata/ml/29-build-an-order-book-model-end-to-end/latency.csv};
\addplot[name path=lo,draw=none,forget plot] table[x=ms,y=lo]{figdata/ml/29-build-an-order-book-model-end-to-end/latency.csv};
\addplot[omProp!25,forget plot] fill between[of=hi and lo];
\addplot[omProp,very thick,mark=*,mark size=1.5pt] table[x=ms,y=pnl]{figdata/ml/29-build-an-order-book-model-end-to-end/latency.csv};
\addplot[omDef,thick,dashed] table[x=ms,y=base]{figdata/ml/29-build-an-order-book-model-end-to-end/latency.csv};
\addplot[black,thin] coordinates {(1e-4,0) (300,0)};
\legend{le modèle ($\pm$ une erreur type),cotation sans modèle}
\end{axis}
\end{tikzpicture}

\omcaption{P\&L net moyen par session de test de vingt minutes en fonction de l'instant de décision, de 150 nanosecondes (modèle compilé) à
150 millisecondes (échelle logarithmique~; zéro donne le même résultat que 150 nanosecondes), avec une bande d'une erreur type sur les six sessions~;
en pointillés~: la même cotation sans modèle. Données~: \texttt{ml\_lobmodel.latency\_table}.}\label{fig:ml:lob:pnl}
\end{omfigure}

Jusqu'à 50 millisecondes rien ne change~: le P\&L reste dans son erreur type de 21.73~\$, et le mark-out à moins de 0.02 ticks de
0.405. Au-delà, l'edge se dégrade rapidement. À 70 millisecondes le P\&L est de 12.40~\$~; à 100 millisecondes il est de 1.60~\$, positif dans
seulement trois sessions~; à 150 millisecondes il est de $-\$12.27$. Par interpolation linéaire sur la grille, le P\&L tombe à zéro vers 106
millisecondes, et à ce que rapporte la cotation sans modèle vers 89 millisecondes~: au-delà, le modèle fait moins bien que l'absence de modèle.

Deux mécanismes produisent la falaise, et la colonne de péremption les sépare. Jusqu'à 70 millisecondes la péremption médiane égale
l'instant de décision~: chaque décision accuse exactement ce retard, et une cotation retirée avec 50 millisecondes de retard est rarement touchée pendant ces 50
millisecondes, car le carnet ne change qu'environ quatre fois par seconde. À 100 millisecondes la péremption médiane est de 154
millisecondes~: dans les rafales les plus intenses, les décisions arrivent plus vite qu'une toutes les 100 millisecondes, $\lambda L$ dépasse un, et elles font la queue
les unes derrière les autres. À 150 millisecondes la péremption médiane est de 802 millisecondes, et l'agent cote sur un carnet vieux de presque une
seconde. Un modèle lent en moyenne l'est encore davantage précisément quand le marché est actif, c'est-à-dire quand ses prévisions comptent.

Le mark-out à une seconde raconte la même histoire du côté des exécutions, qui passe de 0.405 à 0.148 ticks~; il reste positif, car
même un retrait tardif évite certaines des pires exécutions, mais pas assez pour payer le reste.

Pourquoi les microsecondes ne comptent-elles pas ici~? Parce que le temps qui compte se mesure en événements du marché, et non en secondes. Kolm,
Turiel et Westray ont établi que l'horizon effectif de leurs prévisions propres à chaque action valait environ deux variations moyennes de prix. Dans ce
marché, le carnet change environ quatre fois par seconde et personne d'autre ne se dispute les mêmes cotations~: le flux de fond
ne réagit pas à l'agent. Sur une bourse réelle, le haut du carnet d'un titre liquide change bien plus souvent que quatre fois par seconde, et les cotations de l'agent sont
prises à revers par d'autres traders rapides~: Aquilina, Budish et O'Neill, à partir des messages de bourse, observent des courses d'arbitrage de latence environ
une fois par minute et par action du FTSE~100, la course modale durant de 5 à 10 millionièmes de seconde. La même falaise se situerait
d'autant plus près de zéro, et les microsecondes qui n'ont rien rapporté ici seraient toute l'affaire. Le simulateur
répond à la question pour laquelle il a été conçu, le coût de la péremption dans son propre temps des événements~; transposer cette réponse à un marché réel
signifie la recalibrer selon le taux d'événements de ce marché et y ajouter ses concurrents.

\begin{method}[Construire un modèle à court horizon qui négocie]\label{met:ml:lob:method}
\begin{enumerate}
\item Spécifier le marché, la décision, l'action et le budget avant le modèle.
\item Enregistrer les données de recherche avec l'agent de production, afin que les variables reposent sur le même code dans les deux cas~; vérifier que le mode hors ligne donne le même résultat que le mode en ligne.
\item Valider avec des partitions purgées~; choisir le seuil sur des sessions qui n'ont pas servi à l'ajustement du modèle~; tester sur d'autres.
\item Compiler le modèle retenu et vérifier la parité sur des vecteurs partagés~; mesurer sa latence et utiliser la mesure.
\item Négocier le modèle sur une grille d'instants de décision, et présenter le P\&L, les mark-out et la péremption en fonction de l'instant, à côté d'une référence sans modèle.
\item Lire la latence dans le temps d'événements du marché~: l'horizon de l'edge est un nombre de variations de prix, et non de microsecondes.
\end{enumerate}
\end{method}

\section{Tutoriel~: les microsecondes sont de l'argent}\label{tut:ml:build-an-order-book-model-end-to-end:tut}

\begin{tutorial}
\textbf{Objectif.} Exécuter tout le pipeline, compiler le modèle, le négocier à chaque instant de décision, et trouver où s'arrête son edge.
\textbf{État final~:} \cref{tab:ml:lob:latency}, \cref{tab:ml:lob:results}, \cref{fig:ml:lob:pnl}.
\begin{tutsteps}
\item \textbf{Le graphe.}
\omcode{code/firm/lobmodel/firm_lobmodel.py}{344}{358}{Le chapitre comme un seul graphe firm.workflow.}\label{lst:ml:lob:graph}
\item \textbf{Variables à chaque changement du carnet.}
\omcode{code/firm/lobmodel/firm_lobmodel.py}{92}{111}{L'agent enregistre un événement, calcule ses variables en ligne, prévoit et déclare son instant de décision.}\label{lst:ml:lob:onbook}
\item \textbf{Étiquettes.}
\omcode{code/firm/lobmodel/firm_lobmodel.py}{167}{177}{Variables hors ligne et étiquette à horizon fixe en temps d'événements.}\label{lst:ml:lob:data}
\item \textbf{Service en Python.}
\omcode{code/firm/lobmodel/firm_lobmodel.py}{205}{225}{La forêt aplatie parcourue pour tous les arbres à la fois, feuilles ajoutées dans l'ordre.}\label{lst:ml:lob:predict}
\item \textbf{Service en C++20.}
\omcode{code/firm/lobmodel/cpp/lobmodel.hpp}{1}{23}{Le chemin de service C++20~: la boucle de firm.mlinfer et les branches générées.}\label{lst:ml:lob:cpp}
\item \textbf{Service en Rust.}
\omcode{code/firm/lobmodel/rust/src/lib.rs}{5}{23}{Le chemin de service Rust sur le noyau de firm.mlinfer.}\label{lst:ml:lob:rust}
\item \textbf{La décision.}
\omcode{code/firm/lobmodel/firm_lobmodel.py}{113}{130}{Coter chaque côté autorisé au meilleur cours, sauf si la prévision signale qu'il est sur le point d'être pris à revers.}\label{lst:ml:lob:act}
\item \textbf{Exécuter} \texttt{ml\_lobmodel.summary()}, \texttt{latency\_table()}, \texttt{zero\_crossing()}, \texttt{monitors()},
\texttt{fig\_lobmodel.py} (environ trois minutes sur un cœur), et \texttt{bench\_lobmodel.py} une fois pour les latences mesurées.
\end{tutsteps}
\textbf{Pour aller plus loin.} Ajouter au simulateur un second agent de cotation, plus rapide, qui retire ses cotations sur le même signal~: la falaise
devrait se déplacer vers l'écart entre les instants de décision des deux agents.
\end{tutorial}

\section{Construction~: le modèle de carnet d'ordres}\label{bld:ml:build-an-order-book-model-end-to-end:lobmodel}

\begin{build}
\bfield{Rôle} Un modèle à court horizon, depuis les sessions simulées jusqu'à la négociation, avec chaque étape reproductible et son chemin de service dans
trois langages.
\bfield{Interface} \texttt{FEATURES}, \texttt{HORIZON}, \texttt{DEFAULT}~;
\texttt{ModelAgent(predict, threshold, hold\_s, decision\_ns, qty, trade)} (un agent de \texttt{firm.exchsim})~; \texttt{run\_session}, \texttt{record}, \texttt{events\_array},
\texttt{dataset}, \texttt{report}~; \texttt{ForestPredictor}~; les étapes \texttt{stage\_data}, \texttt{stage\_cv},
\texttt{stage\_fit}, \texttt{stage\_tcn}, \texttt{stage\_threshold}, \texttt{stage\_test}~;
\texttt{make\_pipeline(cache\_dir, cfg)}~; \texttt{cpp/lobmodel.hpp} et le crate Rust \texttt{firm\_lobmodel}~; \texttt{make\_lobmodel\_fixture.py}.
\bfield{Règles} La recherche et la production partagent l'agent et le code des variables~; les sessions de rôles différents ne partagent jamais une graine~;
l'instant de décision est une latence mesurée~; les chemins de service concordent bit à bit.
\bfield{Tests d'acceptation} \texttt{code/firm/lobmodel/tests/}~: les variables en ligne sont égales aux variables hors ligne à chaque décision, et les étiquettes
sont égales à la variation du mid~; la forêt plate, la marche NumPy et LightGBM concordent sur l'échantillon de référence, et le pipeline reconstruit la
forêt de l'échantillon~; un instant de décision plus long rend les décisions plus anciennes. \texttt{cpp/lobmodel\_test.cpp} et \texttt{cargo test}~: parité exacte sur
200 vecteurs.
\bfield{Pour aller plus loin} Un agent rapide concurrent~; variables normalisées par l'activité~; le TCN compilé avec le chemin int8 du chapitre~26.
\end{build}

\begin{omsources}
\item P.~N.~Kolm, J.~Turiel et N.~Westray, ``Deep order flow imbalance~: extracting alpha at multiple horizons from the limit
order book'', \emph{Mathematical Finance}, 2023.
\item J.~Sirignano et R.~Cont, ``Universal features of price formation in financial markets~: perspectives from deep learning'',
\emph{Quantitative Finance}, 2019.
\item R.~Cont, A.~Kukanov et S.~Stoikov, ``The price impact of order book events'', \emph{Journal of Financial Econometrics},
2014.
\item Á.~Cartea, S.~Jaimungal et J.~Penalva, \emph{Algorithmic and High-Frequency Trading}, Cambridge University Press, 2015.
\item M.~Aquilina, E.~Budish et P.~O'Neill, ``Quantifying the high-frequency trading `arms race''', \emph{Quarterly Journal of
Economics}, 2022.
\end{omsources}

\section{Exercices}

\begin{exercise}[$\star$]\label{exo:ml:build-an-order-book-model-end-to-end:1}
Pourquoi l'agent enregistre-t-il lui-même ses données de recherche, plutôt que de s'appuyer sur le flux complet du simulateur~?
\end{exercise}

\begin{exercise}[$\star$]\label{exo:ml:build-an-order-book-model-end-to-end:2}
Pourquoi a-t-on abandonné une stratégie preneuse~? Calculer son coût aller-retour en ticks.
\end{exercise}

\begin{exercise}[$\star$]\label{exo:ml:build-an-order-book-model-end-to-end:3}
Le modèle rapporte davantage que la cotation sans lui, mais entre moins souvent. Expliquer à l'aide des mark-outs.
\end{exercise}

\begin{exercise}[$\star\star$]\label{exo:ml:build-an-order-book-model-end-to-end:4}
À 100 millisecondes, la péremption médiane est de 154 millisecondes. Expliquer, avec $\lambda L$, et dire ce qu'un système de production ferait
à ce sujet.
\end{exercise}

\begin{exercise}[$\star\star$]\label{exo:ml:build-an-order-book-model-end-to-end:5}
Le moniteur d'entrée des sessions de test lit jusqu'à 1.60. Le modèle aurait-il dû être arrêté~? Que modifier dans le moniteur~?
\end{exercise}

\begin{exercise}[$\star\star$]\label{exo:ml:build-an-order-book-model-end-to-end:6}
\emph{Trouver la faille.} «~Le modèle compilé est 1\,700 fois plus rapide que l'appel Python de LightGBM, il rapportera donc 1\,700 fois plus
d'argent.~»
\end{exercise}

\begin{exercise}[$\star\star\star$]\label{exo:ml:build-an-order-book-model-end-to-end:7}
Où se situerait le passage par zéro pour une action dont le haut du carnet change 2\,000 fois par seconde, si l'horizon de l'edge est un nombre fixe
de changements du carnet~? Que laisse de côté l'estimation~?
\end{exercise}

\begin{exercise}[$\star\star\star$]\label{exo:ml:build-an-order-book-model-end-to-end:8}
Écrire la spécification de surveillance de ce modèle en production~: les moniteurs, leurs références, leurs seuils et leurs propriétaires.
\end{exercise}

\section{Problème~: les microsecondes valent de l'argent}

\begin{problem}[{Problème du week-end --- les microsecondes valent de l'argent}]\label{pb:ml:build-an-order-book-model-end-to-end:1}
Le pipeline, le modèle et la grille de latences du chapitre.

\medskip\noindent\textbf{Partie I --- Spécification et données.}
\begin{enumerate}
\item Préciser le marché, la décision, l'action et les sessions.
\item Quelles sont les variables, et pourquoi les valeurs hors ligne et en ligne concordent-elles~?
\item Comment l'étiquette est-elle définie, et pourquoi la plupart des étiquettes valent-elles zéro~?
\item Comment le modèle est-il validé, et que choisit la validation croisée~?
\end{enumerate}

\medskip\noindent\textbf{Partie II --- Le modèle.}
\begin{enumerate}[resume]
\item Comment la forêt et le TCN se comparent-ils~?
\item Quel est le coût par prédiction de la forêt compilée, et par rapport à quoi~?
\item Comment la parité est-elle vérifiée~?
\item Comment le seuil est-il choisi~?
\end{enumerate}

\medskip\noindent\textbf{Partie III --- Négociation.}
\begin{enumerate}[resume]
\item Que rapporte le modèle à 150 nanosecondes, et par rapport à quelle référence~?
\item Que montrent les mark-outs~?
\item Que signale le moniteur d'entrée, et pourquoi~?
\item Définir la \omterm{def:ml:build-an-order-book-model-end-to-end:staleness}{péremption de la décision} et expliquer ses deux composantes.
\end{enumerate}

\medskip\noindent\textbf{Partie IV --- Le verdict.}
\begin{enumerate}[resume]
\item Énoncer le \emph{résultat nommé}~: le P\&L net par session et le mark-out à une seconde en fonction de la latence de décision, et la
latence à laquelle l'edge du modèle tombe à zéro.
\item Pourquoi la courbe est-elle plate jusqu'à 50 millisecondes~?
\item Pourquoi chute-t-elle si vite ensuite~?
\item Que changerait-il sur une bourse réelle~?
\item La firme doit-elle investir dans le chemin de service C++ pour ce marché~?
\item Qu'ajouterait-on au simulateur avant de faire confiance à la courbe~?
\item Comment enregistrer et surveiller ce modèle~?
\item En une phrase~: que vaut une microseconde~?
\end{enumerate}
\end{problem}

\section{Questions d'entretien}

\begin{interviewq}[$\star$ \iqroles{chercheur, ingénieur ML}]\label{iq:ml:build-an-order-book-model-end-to-end:1}
Quelles variables utiliseriez-vous pour prévoir le mid de l'action à la seconde suivante~?
\end{interviewq}

\begin{interviewq}[$\star\star$ \iqroles{chercheur}]\label{iq:ml:build-an-order-book-model-end-to-end:2}
Comment validez-vous un modèle dont les étiquettes se chevauchent dans le temps~?
\end{interviewq}

\begin{interviewq}[$\star\star$ \iqroles{ingénieur ML, développeur}]\label{iq:ml:build-an-order-book-model-end-to-end:3}
Comment vous assurez-vous qu'un modèle calcule les mêmes variables en production qu'en recherche~?
\end{interviewq}

\begin{interviewq}[$\star\star$ \iqroles{chercheur, trader}]\label{iq:ml:build-an-order-book-model-end-to-end:4}
Un teneur de marché ajoute une prévision à court horizon. Où sa valeur apparaît-elle, et comment la mesurez-vous~?
\end{interviewq}

\begin{interviewq}[$\star\star$ \iqroles{ingénieur ML, développeur}]\label{iq:ml:build-an-order-book-model-end-to-end:5}
Votre modèle prend 200 microsecondes en Python. Comment le ramenez-vous sous la microseconde, et comment savez-vous qu'il donne toujours les mêmes
réponses~?
\end{interviewq}

\begin{interviewq}[$\star\star\star$ \iqroles{chercheur, développeur}]\label{iq:ml:build-an-order-book-model-end-to-end:6}
Comment décideriez-vous de la somme à consacrer à la latence pour une stratégie donnée~?
\end{interviewq}
